% Korean edition of figs/openq.tex.
\fig[wide]{openq}{질문 5개가 남아 있다. 4개는 테스트로, 1개는 주최측만 답할 수 있다}[미해결 질문, 해결될 때까지 래퍼가 하는 일, 다음 단계]{%
{\tabfont\begin{tabularx}{\linewidth}{>{\RaggedRight\arraybackslash}p{6mm} >{\RaggedRight\arraybackslash}X >{\RaggedRight\arraybackslash}p{50mm} >{\RaggedRight\arraybackslash}p{50mm}}
\thd{} & \thd{질문} & \thd{잠정 동작} & \thd{해결 방법(다음 단계)}\\
Q1 & 클라이언트 코드에서 가져온 \code{CancelOrder}, \code{ModifyOrder}, \code{DeleteOrder} 본문이 동작하는가? 어느 것도 실서버에서 시험하지 않았다 & 래퍼는 쓰지 않는다. 주문은 \code{CancelOrderBatch}로 정지한다 & 웹 페이지가 어떤 주문에 이 기능을 제공할 때 요청을 캡처한다\\\rowsep
Q2 & 취소가 100건 한도에 포함되는가? & 정지한 주문마다 1건으로 센다. 래퍼는 하루 80건으로 제한한다 & 주최측에 문의한다\\\rowsep
Q3 & 비밀번호를 바꾼 뒤에도 저장한 쿠키 묶음이 갱신되는가? & 쿠키 묶음을 7일짜리 자격 증명으로 다룬다(\ref{sec:design}) & 비밀번호를 바꾼 뒤, 바꾸기 전에 저장한 쿠키 묶음으로 갱신해 본다\\\rowsep
Q4 & 서버는 섹터·소형주 위반을 언제 세는가? 장중 아무 때나인가, 장 마감인가? 섹터 한도는 시장 비중을 따라 움직이고, 리밸런싱 중에는 매도가 매수보다 먼저 체결된다는 보장이 없어 비중이 잠시 한도를 넘을 수 있다(\ref{sec:guidelines}) & \code{plan()}은 호출 때마다 피드에서 읽은 한도로 계획 시점의 목표 비중을 검사한다. \code{guidelines()}는 소형주 합계를 보고한다 & 매시간 \code{Violations}와 피드의 섹터 비중을 기록해 비교한다\\\rowsep
Q5 & 주문 상태 \code{Generated}와 \code{Deleted}는 언제 나타나며, 체결되지 않은 채 \code{hm1}에 이른 시간대 주문은 어떻게 끝나는가? & \code{Done}이 아닌 상태는 모두 진행 중으로 본다. \code{plan()}은 남은 크기를 빼고 \code{wait()}은 계속 기다린다. \code{cnclT}가 null인 \code{Done} 주문은 체결로 보고한다 & 서버가 거부한 주문을 포함해 \code{ords}의 모든 상태 변화를 기록한다\\
\end{tabularx}}}{11월 24일에 해결되어 이 목록에서 뺀 것: 실시간 피드가 밀어 보내는지(\ref{sec:feed}), \code{updated}의 형식과 델타 병합의 정확성(\ref{sec:feed}), 동시 갱신(\ref{sec:auth}), \code{Orders} 매개변수(\ref{sec:queries}), 체결·정지된 주문의 최종 상태와 주문 직후 \code{Accepted} 상태(\ref{sec:orders}), 매도 주문 형식(\ref{sec:wrappertests}), \code{qty2ord}의 반올림 방식(내림, \ref{sec:oneshare}), \code{bs_updated}의 내용(\ref{sec:feed}), 순자산의 출처(\ref{sec:oneshare}), 가이드라인의 근거가 되는 피드 필드(\ref{sec:guidelines}).}{DART Quant Team}
